% =============================================================================
\section{Proof of Theorem~\ref{thm:main}}
\label{app:proofs}
% =============================================================================

We restate and prove Theorem~\ref{thm:main} in full.

\begin{theorem}[DPRI Upper Bound on MIA Advantage, restated]
Let $f_\theta$ be trained on $D$ by minimizing empirical risk with an
$L$-Lipschitz loss $\ell$.
Let $P_{\mathrm{in}}$ and $P_{\mathrm{out}}$ be the distributions of
$\ell(f_\theta, x)$ for members and non-members, respectively.
Assume $P_{\mathrm{out}}$ is sub-Gaussian with variance proxy $\sigma^2$.
Then for any sample $x \in D$:
\begin{equation}
  \mathrm{Adv}(\mathcal{A}^*, f_\theta, x)
  \leq
  \frac{L \cdot u(x)}{\sigma \sqrt{2} \cdot \rho(x)^{1/d}},
  \label{eq:bound_app}
\end{equation}
where $u(x) = \min_{x' \neq x} \|x - x'\|_2$ and
$\rho(x)$ is the $k$-NN density estimate of
Equation~\eqref{eq:density}. In the DPRI implementation we use the
monotone surrogate $1/r_k(x)$, which has the same rate under (A3) and
avoids evaluating high-dimensional ball volumes, which underflow for
large $d$.
\end{theorem}

\noindent\textbf{Assumptions.} We make three modeling assumptions. (A1) The
member and non-member loss distributions $P_{\mathrm{in}}$ and
$P_{\mathrm{out}}$ are Gaussian with a common variance $\sigma^2$ and means
$\mu_{\mathrm{in}}, \mu_{\mathrm{out}}$. (A2) The loss $\ell(f_\theta,\cdot)$
is $L$-Lipschitz in its input. (A3) The local nearest-neighbour distance and
density satisfy $r_k(x) \asymp \rho(x)^{-1/d}$, which holds when the density
is approximately constant over the neighbourhood~\cite{biau2015lectures}.

\begin{proof}
\textbf{Step 1 (advantage to KL).}
For the optimal test between $P_{\mathrm{in}}$ and $P_{\mathrm{out}}$,
the advantage equals the total variation distance, which Pinsker's
inequality~\cite{tsybakov2009introduction} bounds by the KL divergence:
\begin{equation}
  \mathrm{Adv}(\mathcal{A}^*) \leq
  \sqrt{\tfrac{1}{2}\,\mathrm{KL}(P_{\mathrm{in}} \| P_{\mathrm{out}})}.
  \label{eq:bh_app}
\end{equation}

\textbf{Step 2 (KL via the loss gap).}
Under (A1), the KL divergence between two Gaussians of common variance is
exact:
\begin{equation}
  \mathrm{KL}(P_{\mathrm{in}} \| P_{\mathrm{out}})
  = \frac{(\mu_{\mathrm{in}} - \mu_{\mathrm{out}})^2}{2\sigma^2}.
\end{equation}
By (A2), the mean losses of a member $x$ and its nearest out-of-sample point
$x^* = \arg\min_{x' \notin D}\|x-x'\|_2$ differ by at most $L\|x-x^*\|_2$.
Therefore
\begin{equation}
  \mathrm{KL}(P_{\mathrm{in}} \| P_{\mathrm{out}})
  \leq \frac{L^2 \|x - x^*\|_2^2}{2\sigma^2}.
  \label{eq:kl_app}
\end{equation}

\textbf{Step 3 (distance via density).}
The member-to-nearest-neighbour distance is $u(x)$ by definition, and by (A3)
the spacing to a fresh draw from $\mathcal{P}$ scales as $\rho(x)^{-1/d}$.
Combining these into the geometric factor $u(x)/\rho(x)^{1/d}$, substituting
into~\eqref{eq:kl_app}, and applying~\eqref{eq:bh_app},
\begin{equation}
  \mathrm{Adv}(\mathcal{A}^*) \leq
  \sqrt{\frac{L^2 u(x)^2}{2\sigma^2 \rho(x)^{2/d}}}
  = \frac{L\, u(x)}{\sigma\sqrt{2}\,\rho(x)^{1/d}}.
\end{equation}
Averaging over $x \in D$ gives the dataset-level bound~\eqref{eq:risk_bound}.
\end{proof}

\noindent\textbf{Remark.} The bound is a modeling result rather than a
worst-case guarantee. It rests on the Gaussian score model (A1) and the local
density-distance relation (A3), both standard but approximate. Its value is
the qualitative prediction it makes, that MIA advantage grows with sample
uniqueness and falls with local density, which Section~\ref{sec:results}
tests directly against measured AUC and finds to hold as a moderate,
statistically significant trend.

\begin{proposition}[Tightness]
\label{prop:tight}
The bound~\eqref{eq:bound_app} is tight up to constants for datasets
drawn from a uniform distribution over a $d$-dimensional hypercube.
\end{proposition}
\begin{proof}
For uniform data, $\rho(x) = c$ (constant) and $u(x) \sim n^{-1/d}$
by standard spacing theory.
The Bayes-optimal MIA advantage for a uniformly distributed dataset
also scales as $O(n^{-1/d})$ (from the Yatracos class bound).
Hence the bound and the true advantage have the same rate.
\end{proof}

Section~\ref{sec:tightness_main} validates
Proposition~\ref{prop:tight} empirically on uniform hypercube data:
the bound factor's fitted decay slopes match the predicted $-1/d$
within $0.03$ for $d \in \{2,4,8\}$.

% =============================================================================
\section{Dataset Details}
\label{app:datasets}
% =============================================================================

The 31 datasets span medical (e.g.\ Breast-W, Pima Diabetes), image (MNIST,
Fashion-MNIST, Digits, OptDigits, PenDigits), network (KDDCup), forest
cover (CovType), credit (German Credit), and a range of tabular sources
(OpenML), together with the canonical MIA benchmarks Purchase100 and
Texas100~\cite{shokri2017membership} and the mobility and recommendation
datasets Gowalla~\cite{cho2011friendship} and
MovieLens~\cite{movielens}.
Datasets larger than 30k samples are subsampled to 30k to keep the
attack-model grid tractable. All datasets are publicly available and loaded
through scikit-learn, OpenML, or their original public sources.
